% Generated draft: requires mathematical and source review.
\documentclass[11pt]{article}
\usepackage[margin=25mm]{geometry}
\usepackage{fontspec}
\setmainfont{Latin Modern Roman}
\usepackage{amsmath,amssymb,mathtools,unicode-math}
\setmathfont{Latin Modern Math}
\usepackage{enumitem,needspace,xurl,hyperref}
\setlength{\parindent}{0pt}
\setlength{\parskip}{5pt}
\setlength{\emergencystretch}{3em}
\newcommand{\N}{\mathbb{N}}
\newcommand{\Z}{\mathbb{Z}}
\newcommand{\Q}{\mathbb{Q}}
\newcommand{\R}{\mathbb{R}}
\newcommand{\C}{\mathbb{C}}
\newcommand{\F}{\mathbb{F}}
\newcommand{\A}{\mathbb{A}}
\newcommand{\PP}{\mathbb{P}}
\newcommand{\E}{\mathbb{E}}
\newcommand{\eps}{\varepsilon}
\newcommand{\dd}{\,\mathrm{d}}
\begin{document}
\section*{UnsolvedMath ID: 700401}
\textbf{Source problem code:} catalogue entry
\subsection{Definitions and mathematical statement}
Let \(G\) be a group. For a finite subgroup \(H\leq G\), define its normalizer by
\[
N_G(H)=\{g\in G:gHg^{-1}=H\}.
\]
Then \(H\trianglelefteq N_G(H)\), so the quotient \(N_G(H)/H\) is defined. Write \(|x|\) for the order of an element \(x\).

The group \(G\) is \emph{biprimitively finite} if, for every finite subgroup \(H\leq G\) and every pair \(x,y\in N_G(H)/H\) whose orders are prime, the subgroup \(\langle x,y\rangle\) is finite. The two prime orders need not be equal.

The group \(G\) is \emph{conjugacy biprimitively finite} if, for every finite subgroup \(H\leq G\) and every pair of conjugate elements \(x,y\in N_G(H)/H\) of prime order, the subgroup \(\langle x,y\rangle\) is finite. Here conjugacy is taken in the quotient \(N_G(H)/H\).

For a fixed prime \(p\), the supplied parenthetical definition of \emph{\(p\)-conjugacy biprimitively finite} is read as requiring the same condition for conjugate elements \(x,y\in N_G(H)/H\) with \(|x|=|y|=p\). This terminology is included in the source definition, but the requested conclusion does not explicitly invoke this variant alone.

A group is \emph{periodic} if every element has finite order. We use \emph{finite rank} in the standard subgroup-generator sense: there exists an integer \(r\geq 0\) such that every finitely generated subgroup of \(G\) can be generated by at most \(r\) elements. A group is \emph{locally finite} if every finitely generated subgroup is finite. An \emph{involution} is an element of order \(2\).

\textbf{Requested statement (V. P. Shunkov, 1978).} Prove that every periodic group of finite rank that is biprimitively finite is locally finite; the parenthetical conjugacy formulation asks for the same conclusion under the hypothesis of conjugacy biprimitively finiteness. Thus the requested implications, with the definitions above, are
\[
\begin{aligned}
&G\text{ periodic, of finite rank, and biprimitively finite}
&&\Longrightarrow G\text{ locally finite},\\
&G\text{ periodic, of finite rank, and conjugacy biprimitively finite}
&&\Longrightarrow G\text{ locally finite}.
\end{aligned}
\]
The source additionally inserts the qualification ``in particular, having no involutions.'' Its intended scope is not explained in the supplied record. We retain the case with the additional condition
\[
\{g\in G:|g|=2\}=\varnothing
\]
as an explicitly mentioned subcase of the stated hypotheses. The supplied wording does not clearly specify any separate conclusion for all periodic groups of finite rank having no involutions without a biprimitively finite hypothesis; no such additional assertion is made here.
\subsection{Short English statement}
Shunkov's problem asks whether a periodic group of finite rank must be locally finite when every pair of prime-order elements in each finite-subgroup normalizer quotient generates a finite subgroup, including the variant that restricts the pairs to conjugate elements. The source specifically mentions groups without involutions, but the scope of that parenthesis requires review. The supplied registry labels the problem open; this label is not an independently verified mathematical assessment.
\subsection{Sources}
\begin{itemize}
\item Supplied local UnsolvedMath registry record: permanent ID \(700401\), exact source problem code \texttt{catalogue entry}, title \emph{Kourovka Notebook Problem 6.59}; upstream revision \texttt{d371ee917ac9647fbb06cd6fc529deec6f3dda82}. The supplied status is \texttt{open}.
\item \emph{The Kourovka Notebook (Unsolved Problems in Group Theory)}, Problem 6.59, proposed by V. P. Shunkov in the 6th issue (1978). The supplied background identifies the 21st issue edition (2026), arXiv identifier \texttt{1401.0300v48}.
\item The supplied collection description attributes editing of Issues 1--21 (1965--2026) to E. I. Khukhro and V. D. Mazurov. These bibliographic details are supplied attribution and have not been independently checked against the cited edition.
\end{itemize}
\end{document}
